\section{Proofs for Verification}
\label{section: verification proofs}

First, we remind useful tools from probability theory.
\begin{lemma}[Hoeffding bound for the binomial distribution]
    \label{lemma:hoeffding:binomial}
    Let $X\sim \mathrm{Binomial}(n,p)$.
    For any $k\le np$,
    \[
    \Pr[X\le k]
    \;\le\;
    \exp\!\left(-\frac{2(np-k)^2}{n}\right),
    \]
    and for any $k\ge np$,
    \[
    \Pr[X\ge k]
    \;\le\;
    \exp\!\left(-\frac{2(np-k)^2}{n}\right).
    \]
    \end{lemma}

    \begin{lemma}[Hoeffding bounds for the hypergeometric distribution]
        \label{lemma:hoeffding:hypergeometric}
        Let $X \sim \mathrm{Hypergeometric}(N,K,n)$ be the number of marked items when drawing $n$ samples without replacement from a population of size $N$ containing $K$ marked elements.
        Then $\mathbb{E}[X]=\frac{K}{N}n$.
        For any $\chi\ge 0$ such that
        \[
        \frac{K}{N}-\chi \ge 0
        \quad\text{and}\quad
        \frac{K}{N}+\chi \le 1,
        \]
        we have
        \[
        \Pr\!\left[
        X \le \Big(\frac{K}{N}-\chi\Big)n
        \right]
        \;\le\;
        \exp(-2\chi^2 n),
        \]
        and
        \[
        \Pr\!\left[
        X \ge \Big(\frac{K}{N}+\chi\Big)n
        \right]
        \;\le\;
        \exp(-2\chi^2 n).
        \]
        \end{lemma}

Now, we prove the main Lemma useful in the security proof of the Verification protocol, namely Lemma~\ref{lemma:security error}, that we hereby re-state.
\securityerror*

\begin{proof}
In the following, we let $\E\in\Pauli_{N\times(n+t)}$ be a fixed Pauli deviation on $N$ rounds, where $N=d+s$, and let $\sigma$ be a partition of $[N]$ drawn uniformly at random.
Then, we can define the following random variables:
\begin{itemize}
    \item $Z$ describing the number of computation rounds affected, meaning the number of $i$ ($\leq d$) yielding a decision bit $\b^{(i)}_1 = y^{(i)}\neq z^\star$ (either by the Server deviation or by inherent failure of the algorithm);
    \item $Y$ describing the number of failed test rounds, meaning the number of $i$ ($>d$) yielding measurement outcomes $\b^{(i)}$ such that $b^{(i)}_{q_i}\neq0$.
\end{itemize}
The statement is thus about bounding the quantity $\max_\E \Pr[Y < w \wedge Z > d/2]$.

\begin{enumerate}[label=\arabic*)]
    \item
    In what follows, we will parametrize the Server cheating strategy by the number of rounds on which it decides to act with a non-trivial deviation.
    We use $\E^{(i)}\in\Pauli_{n+t}$ for the ``per-round Pauli deviation'', meaning such that $\E=\bigotimes_{i\leq N}\E^{(i)}$.
    \item
    With this notation, we can parametrize the deviation on the number of rounds on which it is harmful (see Definition \ref{verification:def:harmful}).
    Let us note $m$ that number. Formally, $$
    m=
    \Big|\left\{ i\leq N \; \text{s.t}\;  W_{\X,\Y}\left( \E^{(i)} \right)\cap Q\neq \emptyset  \right\}\Big|
    .$$ The distinguishing advantage parametrized by $m$ writes
    \begin{equation}
        \label{eq:pd}
        p_d = \max_{m\leq N} \mathrm{Pr}\left[Z > \frac{d}{2} \wedge Y<w\right].
    \end{equation}
    \item
    Now, we focus on the values of $m$ around which the highest values of $p_d$ are met.
    Note that we can always let $m_0\leq N$ and decompose Equation~\ref{eq:pd} into two regimes, $m\leq m_0$ and $m>m_0$.
    \begin{eqnarray}
        \label{eq:pd_dec}
        p_d
        &=& \max
        \left\{
            \max_{0< m\leq m_0}\mathrm{Pr}\left[Z > \frac{d}{2} \wedge Y<w\right],
            \max_{m_0 < m \leq N}\mathrm{Pr}\left[Z > \frac{d}{2} \wedge Y<w\right]
         \right\}
         \\
        &=& \max
        \left\{
            \max_{0< m\leq m_0}\mathrm{Pr}\left[Z > \frac{d}{2}\right] \mathrm{Pr} [Y<w],
            \max_{m_0 < m \leq N}\mathrm{Pr}\left[Z > \frac{d}{2}\right] \mathrm{Pr} [Y<w]
         \right\}
         \\
        &\leq& \max
        \left\{
            \max_{m\leq m_0}\mathrm{Pr}\left[Z > \frac{d}{2}\right] ,
            \max_{m> m_0} \mathrm{Pr} [Y<w]
         \right\}
    \end{eqnarray}
    The reason we did that simplification is the following: when $m$ is upper-bounded, the least probable event is $Z>\frac{d}{2}$, as this event captures the probability of \textit{many computation rounds} to be detected.
    Likewise, when $m$ is lower-bounded, it means that at least $m_0$ rounds are being attacked: in this case the least probable event is to be undetected by those deviations.
    We will in~\ref{item: m_0} develop an intuition on where to place the barrier $m_0$.
    \item
    For the moment, let us notice that in both terms, the maximal value is obtained at $m=m_0$.
    The distinguishing advantage can therefore be written as
    \begin{equation}
        p_d = \max(\epsilon, \nu)
    \end{equation}
    where
    \begin{equation}
        \nu =
        \mathrm{Pr}\left[Z > \frac{d}{2} \;|\; m_0 \; \text{rounds attacked}\right]
        ,\qquad
        \epsilon =
        \mathrm{Pr} [Y<w\;|\; m_0 \; \text{rounds attacked}]
    \end{equation}
    \item With that being said, let us first calculate $\epsilon$.
    First, let $X$ denote a random variable describing the number of test rounds affected by a non-trivial deviation.
    Then, let $x>0$. We can always write
    \begin{eqnarray}
        \epsilon &=&
        \mathrm{Pr} [Y<w\;|\; X>x] \quad \mathrm{Pr}[X>x]
        +
        \mathrm{Pr} [Y<w\;|\; X\leq x] \quad \mathrm{Pr}[X\leq x]
        \\
        &\leq&
        \mathrm{Pr}[X\leq x]
        +
        \mathrm{Pr} [Y<w\;|\; X> x]
    \end{eqnarray}
    From there, we can notice that $X$ is naturally upper-bounded in the stochastic order by a hypergeometrically-distributed random variable $\tilde X \sim H(N, m_0, s)$.
    We use this fact to compute the first term of the sum.
    By using the tail bounds for such the hypergeometric variable, we have for $\chi_\epsilon\geq 0$:
    \begin{equation}
        \mathrm{Pr}\left[
            \tilde X \leq \left(
                \dfrac{m_0}{N} - \chi_\epsilon
             \right)s
         \right] \leq \exp(-2\chi_\epsilon^2 s)
    \end{equation}
    Using this tail bounds, the parametrization on $\chi_\epsilon$ requires to set $x = ({m_0}/{N} - \chi_\epsilon)s$ and we get
    \begin{equation}
        \mathrm{Pr}[X\leq x]
        \leq
        \mathrm{Pr}\left[
            \tilde X \leq x
         \right]
         \leq \exp(-2\chi_\epsilon^2 s)
    \end{equation}
    Now we compute an upper-bound on the second term of the sum, noting that $$\mathrm{Pr} [Y<w\;|\;  X> x] \leq \mathrm{Pr} [Y<w\;|\; X= x]$$ so that $X$ is fixed and parametrized according to $\chi_\epsilon$ as previously.
    Then, we notice that $Y$ conditioned on $X=x$ is upper-bounded in the stochastic order by a binomially-distributed variable $\tilde Y \sim (x, 1/k)$.
    Indeed: the number of affected test rounds has been fixed to $x$, there are $k=|\mathcal Q|$ types of traps, and any non-trivial deviation is detected by at least one type of trap.
    For Protocol~\ref{protocol:verification}, this is guaranteed by Lemma~\ref{lemma:harmful deviations are detected}, where the set of traps $\mathcal Q$ defined by \ref{eq: set of traps} was used. Consequently, $k=1+t$.
    The probability that a given test round detects the non-trivial deviation of the Server is thus at least $1/k$.
    The worst case is when the deviation is detected by only one type of trap, so we can fix the detection rate to $1/k$.
    Using the tail bounds for the binomial variable, we have for
    \begin{equation}
        \mathrm{Pr}\left[
            \tilde Y < w
         \right]
         \leq
         \exp\left[
            -2\dfrac{(x/k-w)^2}{x} \; \right]
    \end{equation}
    for $w < x/k$.
    In the parametrization on $\chi_\epsilon$, this means $\chi_\epsilon < \frac{m_0}{N}-\frac{w}{s}k$.
    In short, we can upper-bound $\epsilon$ as follows:
    \begin{equation}
        \label{eq:epsilon}
        \epsilon \leq
        \min_{\chi_\epsilon\in[0, \frac{m_0}{N}-\frac{w}{s}k]}
        \exp(-2\chi_\epsilon^2 s) + \exp\left[
            -2\dfrac{(
                (\frac{m_0}{N}-\chi_\epsilon)
                \frac{1}{k}-\frac{w}{s}
                 )^2}{
                \frac{m_0}{N}-\chi_\epsilon
            } \; s
          \right]
    \end{equation}

    \item Then, let us calculate $\nu$.
    To this end, let us decompose the random variable $Z$ into $Z_1$ and $Z_2$, where $Z_1$ describes the number of computation rounds affected by a non-trivial deviation while $Z_2$ describes the number of computation rounds non-affected by the Server deviation but that nevertheless yield an incorrect outcome due to the probabilistic nature of the BQP computation.
    Then, for $z_1>0$, $\nu$ can be written as
    \begin{eqnarray}
        \nu &=& \mathrm{Pr}\left[Z_1 + Z_2 > \frac{d}{2}\right] \\
        &\leq& \mathrm{Pr}\left[Z_1 + Z_2 > \frac{d}{2} \;\; | Z_1> z_1\right] \;\; \mathrm{Pr}[Z_1>z_1]
            + \mathrm{Pr}\left[Z_1 + Z_2 > \frac{d}{2} \;\; | Z_1\leq  z_1\right] \;\; \mathrm{Pr}[Z_1\leq z_1] \\
        &\leq& \mathrm{Pr}[Z_1>z_1]
            + \mathrm{Pr}\left[Z_1 + Z_2 > \frac{d}{2} \;\; | Z_1\leq  z_1\right] \\
    \end{eqnarray}
    Let us now study both terms separately.
    The first term can be upper-bounded as follows. $Z_1$ is lower-bounded in the usual stochastic order by a $H(N, m_0, d)$ hypergeometrically distributed random variable $\tilde Z_1$.
    Therefore, using the tails bound for such a distribution, we get for $\chi_\nu>0$, the desired bound if we re-write $z_1 = (m_0/N+\chi_\nu)d$:
    \begin{equation}
        \mathrm{Pr}\left[ Z_1 > \left( \dfrac{m_0}{N} + \chi_\nu \right)d \right] \leq \exp(-2\chi_\nu^2 d)\; .
    \end{equation}
    For the first term, let us note that
    \begin{equation}
        \mathrm{Pr}\left[Z_1 + Z_2 > \frac{d}{2} \;\; | Z_1\leq  z_1\right] \leq \mathrm{Pr}\left[Z_2 > \frac{d}{2} - z_1 \;\; | Z_1\leq  z_1\right] \leq \mathrm{Pr}\left[Z_2 > \frac{d}{2} - z_1 \;\; | Z_1= z_1\right]
    \end{equation}
    Now, with $Z_1$ fixed to $z_1$, $Z_2$ becomes upper-bounded in the usual stochastic order by a $B(d-z_1, \bqp)$ binomially-distributed random variable $\tilde Z_2$, as among the $d$ computation rounds, only $d-z_1$ are non-affected by a deviation, and the failure probability is $\bqp$ for each.
    The tails bound for $X \sim B(n,\bqp)$ states that, for $n\bqp \le k$
    \begin{equation}
        \mathrm{Pr}[X> k] \leq \exp\left[
            -2\dfrac{(n\bqp-k)^2}{n}
         \right]
    \end{equation}
    Here, replacing $n$ by $d-z_1 = d(1-\frac{m_0}{N}-\chi_\nu)$, and $k$ by $\frac{d}{2}-z_1 = d(\frac{1}{2}-\frac{m_0}{N}-\chi_\nu)$, we get for $(d-z_1)\bqp < \frac{d}{2}-z_1$, or equivalently $\chi_\nu < \frac{1-2\bqp}{2-2\bqp}-\frac{m_0}{N}$ :
    \begin{equation}
        \mathrm{Pr}\left[Z_2 > \frac{d}{2}-z_1\right] \leq \exp\left[
            -2
            \dfrac{
                [\left( 1-\frac{m_0}{N}-\chi_\nu \right)d\bqp - \left( \frac{1}{2}-\frac{m_0}{N}-\chi_\nu \right)d]^2
            }{
                \left( 1-\frac{m_0}{N}-\chi_\nu \right)d
            }
        \right]
    \end{equation}
    Finally, after simplifying and putting back the pieces together, we get the following
    \begin{equation}
        \label{eq: nu}
        \nu \leq \min_{
            \chi_\nu \in [0, \frac{1-2\bqp}{2-2\bqp}-\frac{m_0}{N}]
        }
        \exp[ -2\chi_\nu^2 d ] + \exp\left[
            -2
            \dfrac{
                \left( (1-\frac{m_0}{N}-\chi_\nu)(1-\bqp)-\frac{1}{2} \right)^2
            }{1-\frac{m_0}{N}-\chi_\nu} d
         \right]
    \end{equation}

    \item \label{item: m_0}
    Lastly, we need to develop an intuition on where to place $m_0$.
    For this barrier to be meaningful, let the following scenario, in which a fraction $f$ of the total rounds are attacked by the Server with a non-trivial deviation.
    Note that to obtain the classical output, the Client performs a majority vote over the obtained results.
    Therefore, we are interested in quantifying the critical fraction that corrupts more than half of the outcomes.
    However, we know that there is an inherent probability $\bqp$ of failure for a BQP computation.
    Hence, our criteria becomes the following: if $f$ is the fraction of rounds attacked by a Server deviation, then this must satisfy $(1-\bqp)(1-f)>1/2$, in the sense the fraction $1-f$ of non-attacked rounds that succeed (with probability $1-\bqp$) are more than one half of the total rounds.
    This becomes $f<(1-2\bqp)/(2-2\bqp)$. We then let $\varphi>0$ and parametrize $m_0$ according to how much the Server is far or close to attacking the critical fraction.
    \begin{equation}
        m_0 = \left( \dfrac{1-2\bqp}{2-2\bqp}-\varphi\right)N
    \end{equation}
    We can thus replace $m_0$ by this expression parametrized by $\varphi$ in Equations~\ref{eq:epsilon} and~\ref{eq: nu}.
    \item Then, it means that for different values of $\varphi$, different values of $\epsilon, \nu$ are obtained: these have become function of $\varphi$.
    \begin{equation}
        \label{eq: epsilon_varphi}
        \epsilon(\varphi) \leq
        \min_{\chi_\epsilon\in[0, (\frac{1-2\bqp}{2-2\bqp}-\frac{w}{s}k)-\varphi ]}
        \exp(-2\chi_\epsilon^2 s) + \exp\left[
            -2\dfrac{( (\frac{m_0}{N}-\chi_\epsilon) \frac{1}{k}-\frac{w}{s} )^2}{ \frac{m_0}{N}-\chi_\epsilon } \; s
          \right]
    \end{equation}

    \begin{equation}
        \label{eq: nu_varphi}
        \nu(\varphi) \leq \min_{ \chi_\nu \in [0, \varphi] }
        \exp[ -2\chi_\nu^2 d ] + \exp\left[
            -2 \dfrac{ \left( (1-\frac{m_0}{N}-\chi_\nu)(1-\bqp)-\frac{1}{2} \right)^2 }{1-\frac{m_0}{N}-\chi_\nu} d
         \right]
    \end{equation}

    However, the distinguishing advantage can always be upper-bounded by the minimal one.
    Furthermore, note that for the bounds on $\chi_\epsilon$ in Equation~\ref{eq: epsilon_varphi} to be well-defined, we need to have $\varphi < \frac{1-2\bqp}{2-2\bqp}-\frac{w}{s}k$.
    Finally, the distinguishing advantage can be written as follows:
    \begin{equation}
        p_d \leq \min_{\varphi \in [0, \frac{1-2\bqp}{2-2\bqp}-\frac{w}{s}k]} (\max(\nu(\varphi), \epsilon(\varphi)))
    \end{equation}
    where again, as a reminder, $k=|\mathcal Q|$.
    For Protocol \ref{protocol:verification}, $k=1+t$.

    \item Now, to set a concrete bound, we can set the buffers $\chi_\epsilon, \chi_\nu, \varphi$ to concrete values.
    Indeed, one can re-define the following constants:
    \begin{itemize}
        \item $\alpha = (1-2\bqp)/(2-2\bqp)$
        \item $\Delta = \alpha-wk/s$
        \item $m_0/N = \alpha-\varphi$
    \end{itemize}
    With this, the range of $\varphi$ is $[0, \Delta]$, so we can set $\varphi = \Delta/2$, in the middle of the range.
    Then, the range for $\chi_\epsilon$ becomes $[0, \Delta/2]$, thus we can set $\chi_\epsilon = \chi_0:=\Delta/4$, again in the middle of the range.
    Finally, for $\chi_\nu$ in the expression of $\nu$, the range is $[0, \varphi]=[0, \Delta/2]$ so we can again set $\chi_\nu=\chi_0$.
    When we replace these buffer values in the above expressions, we get values of $\epsilon, \nu$ that upper-bound the actual values suggested by the minimisation problem above, but nevertheless negligible in the number of rounds $d, s$:
    \begin{align}
        \epsilon &\leq \exp( -2(\Delta/4)^2 s ) + \exp\left( -2\dfrac{ ((\alpha-3\Delta/4)/k-\frac{w}{s})^2 }{\alpha-3\Delta/4} s \right)
        \\
        \nu&\leq \exp( -2(\Delta/4)^2 d ) + \exp\left( -2\dfrac{ (((1-\alpha)+\Delta/4)(1-\bqp)-1/2)^2 }{(1-\alpha)+\Delta/4} d \right)
    \end{align}

\end{enumerate}
\end{proof}